% ============================================================
%  Глава 2 — Синтез регуляторов: модальный + ЛКР
% ============================================================
\chapter{Синтез регуляторов: модальный и оптимальный}
\label{ch:synthesis}

В этой главе строится статический регулятор $\vect{u} = -\mat{K}\,\vect{x}$,
размещающий полюса замкнутой системы в желаемых точках.
Излагаются два подхода (по разделам 3.1 и 3.3 учебника):
\textbf{модальный} (выбор полюсов вручную) и
\textbf{оптимальный ЛКР} (минимизация интегрального критерия).

\section{Постановка задачи стабилизации}
\label{sec:stabilization}

Дана линейная стационарная модель (по результатам главы~\ref{ch:model}):
\begin{equation}
  \dot{\vect{x}}(t) = \mat{A}\,\vect{x}(t) + \mat{B}\,\vect{u}(t),
  \qquad \vect{x}(0) = \vect{x}^0.
  \label{eq:plant_continuous}
\end{equation}

Требуется построить регулятор статической обратной связи:
\begin{equation}
  \boxed{\vect{u}(t) = -\mat{K}\,\vect{x}(t),}
  \label{eq:state_feedback}
\end{equation}
такой, что замкнутая система $\dot{\vect{x}} = (\mat{A}-\mat{B}\mat{K})\vect{x}$
асимптотически устойчива, и переходный процесс удовлетворяет требованиям
по длительности и перерегулированию.

\section{Модальный регулятор (pole placement)}
\label{sec:modal}

\subsection{Идея метода}

Зададим желаемые собственные значения замкнутой системы
$\{\lambda_i^{\text{уст}}\}_{i=1}^n$ и подберём $\mat{K}$ так,
чтобы спектр $\sigma(\mat{A}-\mat{B}\mat{K})$ совпал с этим множеством.

Согласно п.~3.1.3 учебника, для гарантии экспоненциального переходного
процесса заданной длительности $t_n$ выбираются вещественные собственные
значения, удовлетворяющие
\begin{equation}
  \min_i |\lambda_i^{\text{уст}}| \;\geq\; \frac{t_n}{3}.
  \label{eq:eigval_bound}
\end{equation}

\subsection{Каноническая форма Фробениуса}

Для одномерного входа ($r=1$) модальный регулятор синтезируется через
приведение к канонической форме (п.~3.1.2 учебника):
\begin{equation}
  \mat{A}_f = \begin{pmatrix}
    0 & 1 & 0 & \cdots & 0\\
    0 & 0 & 1 & \cdots & 0\\
    \vdots&&&\ddots&\vdots\\
    -a_n & -a_{n-1} & \cdots & & -a_1
  \end{pmatrix},
  \quad
  \vect{b}_f = \begin{pmatrix}0\\0\\ \vdots \\0\\1\end{pmatrix}.
  \label{eq:frobenius_form}
\end{equation}

Преобразование подобия $\mat{Q} = \mat{S}_y\mat{T}$, где
\begin{equation}
  \mat{T} = \begin{pmatrix}
    a_{n-1} & a_{n-2} & \cdots & a_1 & 1\\
    a_{n-2} & a_{n-3} & \cdots & 1 & 0\\
    \vdots & & & & \vdots \\
    1 & 0 & \cdots & & 0
  \end{pmatrix},
  \label{eq:T_matrix}
\end{equation}
$a_i$ --- коэффициенты характеристического полинома исходной матрицы
$\chi_n(\lambda) = \det(\lambda \mat{I} - \mat{A}) = \lambda^n + a_1\lambda^{n-1}+\ldots+a_n$.

Параметры регулятора в исходном базисе (формула 3.4 учебника):
\begin{equation}
  \boxed{\mat{K} = \big(\mat{Q}^{\T}\big)^{-1}(\vect{a} - \vect{a}_{\text{уст}}),}
  \label{eq:modal_gain}
\end{equation}
где $\vect{a}_{\text{уст}} = (a_n^{\text{уст}}, a_{n-1}^{\text{уст}}, \ldots, a_1^{\text{уст}})\T$ ---
коэффициенты желаемого характеристического полинома
$\prod_{i=1}^n(\lambda - \lambda_i^{\text{уст}})$.

\subsection{Многомерный вход}

При $r > 1$ (как у нашего робота: $u_v$ и $u_\omega$) канонизация по
\eqref{eq:frobenius_form} напрямую неприменима. На практике используют
функцию \texttt{place(A,B,p)} пакета Control System Toolbox, реализующую
алгоритм Tits--Yang (численно устойчивый). Для нашей задачи в MATLAB:
\begin{filebox}[title={\filepath{matlab/design\_modal.m}}]
\begin{lstlisting}[style=matlabstyle]
poles_des = [-2; -3; -4; -8; -10];   % желаемые полюса
K_modal   = place(Ad, Bd, exp(poles_des * Ts));   % дискретный
\end{lstlisting}
\end{filebox}

\section{Оптимальный регулятор (ЛКР)}
\label{sec:lqr}

Подход модального синтеза требует ручного выбора полюсов, что
не масштабируется при многомерном состоянии. \textbf{Линейно-квадратичный
регулятор} (ЛКР) автоматически минимизирует интегральный критерий качества
(п.~3.3 учебника):

\begin{formulabox}[title={Критерий качества}]
\begin{equation}
  J = \int_0^\infty \big(\,\vect{x}\T\,\mat{Q}\,\vect{x} + \vect{u}\T\,\mat{R}\,\vect{u}\,\big)\,dt,
  \qquad \mat{Q} = \mat{Q}\T \succeq 0, \quad \mat{R} = \mat{R}\T \succ 0.
  \label{eq:lqr_cost}
\end{equation}
\end{formulabox}

Здесь $\mat{Q}$ штрафует отклонение состояния (точность),
$\mat{R}$ --- расход управления (энергию). Подбор соотношения $\mat{Q}/\mat{R}$
определяет компромисс «точность $\leftrightarrow$ затраты».

\subsection{Уравнение Риккати}

Оптимальное управление (формула 3.10 учебника) имеет вид
\begin{equation}
  \vect{u}^*(t) = -\mat{R}^{-1}\mat{B}\T\mat{P}\,\vect{x}(t) = -\mat{K}\,\vect{x}(t),
  \label{eq:lqr_optimal}
\end{equation}
где $\mat{P}$ --- симметричная положительно определённая матрица,
удовлетворяющая алгебраическому уравнению Риккати (CARE):

\begin{formulabox}[title={Непрерывное уравнение Риккати}]
\begin{equation}
  \mat{A}\T\mat{P} + \mat{P}\mat{A} - \mat{P}\mat{B}\mat{R}^{-1}\mat{B}\T\mat{P} = -\mat{Q}.
  \label{eq:care}
\end{equation}
\end{formulabox}

\textbf{Дискретный аналог} (DARE):
\begin{equation}
  \mat{A}\T\mat{P}\mat{A} - \mat{P} - \mat{A}\T\mat{P}\mat{B}\bar{\mat{R}}^{-1}\mat{B}\T\mat{P}\mat{A} = -\mat{Q},
  \quad \bar{\mat{R}} = \mat{R} + \mat{B}\T\mat{P}\mat{B},
  \label{eq:dare}
\end{equation}
матрица регулятора $\mat{K}_d = \bar{\mat{R}}^{-1}\mat{B}\T\mat{P}\mat{A}_d$.

\subsection{Условия существования решения}

По п.~3.3 учебника, единственное положительно определённое $\mat{P}$
существует при выполнении трёх условий:
\begin{enumerate}
  \item Объект $(\mat{A},\mat{B})$ управляем
        ($\operatorname{rank}\mat{S}_y = n$);
  \item $\mat{R} \succ 0$;
  \item $\mat{Q} \succ 0$, либо $\mat{Q} = \mat{H}\T\mat{H}$ при наблюдаемой паре $(\mat{A}\T, \mat{H}\T)$.
\end{enumerate}

Для робота все три условия выполнены при $v_0 > 0$.

\subsection{Выбор весовых матриц}

Согласно рекомендации [13] из учебника (п.~3.3 после формулы 3.12),
естественный выбор:
\begin{equation}
  \mat{R} = \mat{I}_r,
  \qquad
  \tilde{\mat{Q}} = \operatorname{diag}\!\left(\tfrac{1}{(y^{\text{ст}})^2},\,\tfrac{1}{(t_{\max}/3)^2}\right),
  \label{eq:weight_choice}
\end{equation}
где $y^{\text{ст}}$ --- допустимая статическая ошибка, $t_{\max}$ ---
требуемая длительность переходного процесса.

\textbf{Для робота SAMURAI} разумные параметры:
\begin{equation}
  \mat{Q} = \operatorname{diag}(10,\,10,\,5,\,1,\,1),
  \qquad
  \mat{R} = \operatorname{diag}(1,\,1).
  \label{eq:Q_R_robot}
\end{equation}
Большие веса на $p_x, p_y, \theta$ означают приоритет точности позиции
и курса; малые веса на $v, \omega$ --- допускаем заметные отклонения
скоростей. $\mat{R} = \mat{I}$ --- штрафуем затраты управления равномерно.

\begin{notebox}
\textbf{«Сделать поведение робота ошибочным»} в терминах ЛКР означает:
\begin{itemize}
  \item Заменить $\mat{Q}$ на отрицательно определённую --- алгоритм
        не сходится или даст нестабильный регулятор;
  \item Поменять знак $\mat{R}$ или сделать $\mat{R}$ очень малой ---
        получим избыточно агрессивное управление с насыщением и
        автоколебаниями;
  \item Намеренно использовать желаемые полюса в правой полуплоскости
        $\lambda^{\text{уст}}_i > 0$ для модального регулятора ---
        замкнутая система будет неустойчива.
\end{itemize}
Эти эксперименты безопасны в симуляторе, но опасны на реальном железе ---
проводить только с поднятыми колёсами.
\end{notebox}

\section{Регулирование с ненулевой уставкой}
\label{sec:setpoint}

ЛКР стабилизирует систему в нулевой точке. Для отслеживания
заданной траектории $\vect{x}_0 \neq \vect{0}$ (например, цели Nav2)
требуется дополнительное управляющее воздействие (п.~3.2 учебника,
формула 3.7):
\begin{equation}
  \vect{u}(t) = -\mat{K}\,\vect{x}(t) + \vect{u}_{\text{уст}}(t),
  \quad
  \vect{u}_{\text{уст}} = -\big(\mat{C}\mat{A}_z^{-1}\mat{B}\big)^{-1}\,\vect{y}_{\text{уст}}(t),
  \label{eq:setpoint_feedforward}
\end{equation}
где $\mat{A}_z = \mat{A} - \mat{B}\mat{K}$ --- матрица замкнутой системы.
Это \textbf{прямая связь по уставке} (feedforward), компенсирующая статическую
ошибку при ступенчатом изменении задания.

Для робота уставка $\vect{y}_{\text{уст}} = (p_x^*, p_y^*, \theta^*, v^*, 0)\T$
формируется из заданной точки и желаемой скорости подхода.

\section{Наблюдатель состояния (Луенбергер)}
\label{sec:observer}

Регулятор \eqref{eq:state_feedback} использует полный вектор $\vect{x}(t)$,
но реальные датчики дают только $\vect{y} = \mat{C}\vect{x}$. Согласно
п.~3.4.1 учебника, состояние восстанавливается полноразмерным наблюдателем:

\begin{formulabox}[title={Уравнение наблюдателя (Luenberger)}]
\begin{equation}
  \dot{\hat{\vect{x}}}(t) = \mat{A}\,\hat{\vect{x}} + \mat{B}\,\vect{u}
       + \mat{L}\big(\vect{y} - \mat{C}\hat{\vect{x}}\big).
  \label{eq:observer}
\end{equation}
\end{formulabox}

Ошибка восстановления $\vect{e} = \vect{x} - \hat{\vect{x}}$ удовлетворяет:
\begin{equation}
  \dot{\vect{e}} = (\mat{A} - \mat{L}\mat{C})\,\vect{e}.
  \label{eq:observer_error}
\end{equation}

Матрица $\mat{L}$ выбирается так, чтобы $(\mat{A}-\mat{L}\mat{C})$
имела заданные собственные значения, обеспечивающие сходимость
$\vect{e}(t) \to 0$ существенно быстрее динамики регулятора.
Стандартная рекомендация (п.~3.4.2 учебника): полюса наблюдателя в $5\ldots10$ раз
быстрее полюсов регулятора:
\begin{equation}
  \tilde{\lambda}^{\text{набл}}_i = 5 \cdot \lambda^{\text{уст}}_i.
  \label{eq:observer_eigvals}
\end{equation}

\subsection{Принцип разделения}

При совместном использовании регулятора и наблюдателя замкнутая
система имеет блочную матрицу (формула 3.23 учебника):
\begin{equation}
  \mat{A}_c = \begin{pmatrix}
    \mat{A}+\mat{B}\mat{K} & -\mat{B}\mat{K}\\
    \mat{0} & \mat{A}-\mat{L}\mat{C}
  \end{pmatrix}.
  \label{eq:separation}
\end{equation}
Спектр $\sigma(\mat{A}_c) = \sigma(\mat{A}+\mat{B}\mat{K}) \cup \sigma(\mat{A}-\mat{L}\mat{C})$ ---
полюса регулятора и наблюдателя независимы (\textbf{принцип разделения}).
Это позволяет проектировать $\mat{K}$ и $\mat{L}$ \textbf{независимо}.

\subsection{Условие наблюдаемости}

Существование $\mat{L}$, обеспечивающего любые желаемые собственные
значения, требует полной наблюдаемости (п.~3.4.1):
\begin{equation}
  \operatorname{rank}\mat{S}_n = n,
  \quad
  \mat{S}_n = \begin{pmatrix} \mat{C}\T & \mat{A}\T\mat{C}\T & \cdots & (\mat{A}\T)^{n-1}\mat{C}\T \end{pmatrix}.
  \label{eq:observability}
\end{equation}
При $\mat{C} = \mat{I}$ (как в нашей модели --- все 5 координат
наблюдаемы) это условие выполнено тривиально.

\section{Численный пример для робота}
\label{sec:numerical_example}

Применим теорию к модели \eqref{eq:AB_numerical} при $v_0 = 0.2$ м/с,
$T_s = 0.05$ с. Дискретизация ZOH даёт:
\begin{equation*}
  \mat{A}_d \approx \begin{pmatrix}
    1 & 0 & 0 & 0.0488 & 0\\
    0 & 1 & 0.01 & 0 & 0.000244\\
    0 & 0 & 1 & 0 & 0.0488\\
    0 & 0 & 0 & 0.717 & 0\\
    0 & 0 & 0 & 0 & 0.607
  \end{pmatrix}.
\end{equation*}

Решение DARE с $\mat{Q} = \operatorname{diag}(10,10,5,1,1)$, $\mat{R} = \mat{I}_2$
(вычисление в \filepath{matlab/design\_lqr.m}, см. главу~\ref{ch:implementation})
даёт регулятор:
\begin{equation}
  \mat{K}_{\text{LQR}} \approx \begin{pmatrix}
    2.84 & 0 & 0 & 1.21 & 0\\
    0 & 0.95 & 1.62 & 0 & 0.74
  \end{pmatrix}.
  \label{eq:K_lqr_robot}
\end{equation}
Структура $\mat{K}$ отражает связи: $u_v$ зависит от $p_x$ и $v$
(продольный канал), $u_\omega$ --- от $p_y, \theta, \omega$
(боковой канал, через рыскание).

Собственные значения замкнутой непрерывной системы $\mat{A} - \mat{B}\mat{K}$
лежат в левой полуплоскости (асимптотическая устойчивость), модули собственных
значений $|\lambda_i|$ соответствуют постоянным времени $0{,}3\ldots 1$ с,
что обеспечивает время регулирования $\sim 1{,}5$ с.

\begin{notebox}
ЛКР даёт автоматический компромисс «точность/затраты» через выбор
$\mat{Q},\mat{R}$. Для расширения на \textbf{ограниченные} управления
(физические лимиты: $|u_v| \leq 0{,}3$ м/с, $|u_\omega| \leq 2$ рад/с)
переходим к MPC --- следующая глава.
\end{notebox}
